% ARQUIVO GERADO AUTOMATICAMENTE por article/scripts/gerar_resultados.py
% Fonte: data/outputs/*.parquet — NÃO EDITAR À MÃO.
\begin{table}[htbp]
\centering
\caption{Tempo de execução da interface por interação (segundos), sem restrição e com 2 núcleos}
\label{tab:app_benchmark}
\resizebox{\textwidth}{!}{%
\begin{tabular}{llccccccc}
\toprule
 & & \multicolumn{3}{c}{\textbf{12 núcleos lógicos}} & \multicolumn{3}{c}{\textbf{2 núcleos (afinidade restrita)}} & \\
\cmidrule(lr){3-5}\cmidrule(lr){6-8}
\textbf{Edição} & \textbf{Atributo} & \textbf{1.ª carga} & \textbf{Inter. (média)} & \textbf{Inter. (máx.)} & \textbf{1.ª carga} & \textbf{Inter. (média)} & \textbf{Inter. (máx.)} & \textbf{Parquets (MB)} \\
\midrule
2019 & Renda (Q006) & 0,18 & 0,145 & 0,162 & 0,17 & 0,120 & 0,121 & 0,84 \\
2019 & Raça/cor & 0,26 & 0,131 & 0,139 & 0,23 & 0,118 & 0,127 & 0,84 \\
2020 & Renda (Q006) & 0,18 & 0,145 & 0,183 & 0,17 & 0,121 & 0,124 & 0,87 \\
2020 & Raça/cor & 0,18 & 0,155 & 0,254 & 0,16 & 0,137 & 0,220 & 0,87 \\
2021 & Renda (Q006) & 0,20 & 0,182 & 0,320 & 0,18 & 0,157 & 0,272 & 0,83 \\
2021 & Raça/cor & 0,18 & 0,135 & 0,138 & 0,16 & 0,117 & 0,119 & 0,83 \\
\bottomrule
\end{tabular}}
\fonte{Elaborado pelo autor a partir de \texttt{app\_benchmark.parquet}. 72 de 72 execuções abaixo de 2~s. Ambiente local: Windows 11, Python 3.12.4, Streamlit 1.65.0. A restrição a 2 núcleos aproxima, mas não reproduz, o ambiente de 2 vCPUs do Hugging Face Spaces.}
\end{table}
